% Generado por Experimentos/generar_resultados_latex.py desde extraido.
% No editar a mano: volver a ejecutar el script.

\begin{table}[ht]
\centering
\footnotesize
\begin{tabular}{c c c c}
\hline
Ciclo & $\alpha=0$ & $\alpha=0.5$ & $\alpha=1$ \\
\hline
0 ($M_0$) & \multicolumn{3}{c}{$30.46$ [$30.40$; $30.51$]} \\
1 & $30.49$ [$30.46$; $30.52$] & $31.61$ [$31.45$; $31.76$] & $32.92$ [$32.81$; $33.00$] \\
2 & $30.46$ [$30.32$; $30.57$] & $31.87$ [$31.57$; $32.03$] & $34.46$ [$34.33$; $34.73$] \\
3 & $30.47$ [$30.43$; $30.51$] & $31.86$ [$31.75$; $32.06$] & $36.11$ [$35.74$; $36.32$] \\
4 & $30.40$ [$30.38$; $30.42$] & $31.90$ [$31.85$; $31.96$] & $38.03$ [$37.73$; $38.53$] \\
5 & $30.47$ [$30.40$; $30.51$] & $31.90$ [$31.69$; $32.18$] & $39.30$ [$39.01$; $39.76$] \\
\hline
\end{tabular}
\caption{Perplejidad sobre referencias originales: media y rango [mínimo; máximo] de las tres cadenas por ciclo y proporción sintética $\alpha$. $M_0$ es compartido por las tres proporciones dentro de cada repetición.}
\label{tab:medias_human_perplexity}
\end{table}
